\documentclass[twocolumn,11pt,a4paper]{article}
\usepackage[utf8]{inputenc}
\usepackage[french]{babel}
\usepackage[T1]{fontenc}
\usepackage{amsmath, amssymb}
\usepackage{graphicx}
\usepackage{geometry}
\geometry{margin=1.5cm}
\usepackage[none]{hyphenat}
\usepackage{booktabs}
\usepackage{float}
\usepackage{caption}
\captionsetup{skip=3pt, font=small}

\title{Minimisation De La Distance Maximum Entre Les Particules}
\author{Groupe 22 : Othmane CHAOUI || Omar DAHMAN || Rayen JALOUALI}
\date{23 avril 2026}


\begin{document}

\maketitle

\section{Critère}

Le critère de notre groupe 22 est de minimiser la distance maximale entre les particules lors de la première itération après T secondes.

\section{Description de l'innovation}

Notre innovation consiste à regrouper ou rassembler les particules vers le centre de notre environnement, ce qui entraîne la minimisation de la distance entre les particules. On identifie deux cas principaux :\\
\begin{itemize}
    \item \textbf{T = 0 :} c'est-à-dire la première itération, donc on met les particules au centre de l'environnement avec un vecteur vitesse nul.
    \item \textbf{T > 0 :} on change le vecteur vitesse des particules tel qu'elles se dirigent vers le centre de l'environnement.
\end{itemize}

\section{Protocole expérimental}

Pour évaluer notre méthode, on a généré \textbf{100 simulations} avec des paramètres choisis aléatoirement. Pour chaque simulation, on compare deux versions : une \textbf{sans notre méthode} (les particules partent vers la droite par défaut) et une \textbf{avec notre méthode}. Les deux simulations tournent pendant $T$ secondes, puis on mesure la distance maximale entre deux particules.

Les paramètres sont tirés au hasard dans les intervalles suivants :

\begin{table}[H]
\centering
\caption{Plages des paramètres utilisés}
\small
\begin{tabular}{ll}
\toprule
Paramètre & Valeurs \\
\midrule
$n$ (nombre de particules) & 2 à 50 \\
$w$ (largeur)              & 10 à 200 \\
$h$ (hauteur)              & 10 à 200 \\
$r$ (rayon de répulsion)   & 1 à 20 \\
$dt$ (pas de temps)        & 0.1 à 1.0 \\
$T$ (durée en secondes)    & 0 à 10 \\
\bottomrule
\end{tabular}
\end{table}

On mesure trois choses sur chaque simulation :
\begin{enumerate}
    \item La \textbf{distance maximale} entre deux particules à la fin (le critère à minimiser).
    \item Le \textbf{temps de calcul} de notre méthode en millisecondes.
    \item Le \textbf{rapport d'amélioration} : distance sans méthode divisée par distance avec méthode. Un rapport supérieur à 1 signifie que notre méthode a réduit la distance maximale.
\end{enumerate}

\section{Résultats}

\subsection{Effet sur la distance maximale}

\begin{figure}[H]
  \centering
  \includegraphics[width=\linewidth]{results/fig1_crit_vs_n.pdf}
  \caption{Distance maximale selon $n$}
\end{figure}

\begin{figure}[H]
  \centering
  \includegraphics[width=\linewidth]{results/fig2_crit_vs_T.pdf}
  \caption{Distance maximale selon $T$}
\end{figure}

On voit que plus la simulation dure longtemps (grand $T$), plus les particules s'éloignent les unes des autres. La différence entre les deux courbes reste faible : notre méthode n'apporte pas un grand avantage pour les longues simulations.

\begin{table}[H]
\centering
\caption{Résultats moyens par tranche de nombre de particules}
\scriptsize
\begin{tabular}{@{}ccccc@{}}
\toprule
Nb part. & Nb sim. & Dist. opt. & Dist. base & Rapport \\
\midrule
2 à 10  & 17 & 6.01 & 6.11 & 1.35 \\
11 à 20 & 22 & 6.19 & 6.73 & 0.82 \\
21 à 30 & 19 & 7.40 & 7.63 & 0.78 \\
31 à 40 & 26 & 7.43 & 6.04 & 0.90 \\
41 à 50 & 16 & 9.51 & 8.44 & 1.15 \\
\bottomrule
\end{tabular}
\end{table}

Avec peu de particules (2 à 10), le rapport est de 1.35 : notre méthode fonctionne bien. Avec plus de particules (11 à 40), le rapport passe sous 1, ce qui veut dire que notre méthode fait parfois moins bien. Cela s'explique par les répulsions : quand beaucoup de particules convergent vers le même centre, elles se repoussent entre elles et repartent dans toutes les directions.

\subsection{Rapport d'amélioration selon la durée $T$}

\begin{figure}[H]
  \centering
  \includegraphics[width=\linewidth]{results/fig3_ratio_vs_T.pdf}
  \caption{Rapport d'amélioration selon $T$ -- au-dessus de la ligne rouge = amélioration}
\end{figure}

Les points sont éparpillés autour de 1 quelle que soit la durée $T$. Notre méthode améliore les résultats dans \textbf{46 simulations sur 100} et les dégrade dans les 54 autres. Le rapport moyen global est de 0.975.

\subsection{Temps de calcul}

Notre méthode est très rapide : en moyenne 0.1 ms, au maximum 10 ms. Elle ne ralentit pas le programme.

\section{Conclusion}

Notre méthode de convergence vers le centre est simple et très rapide. Elle donne de bons résultats quand il y a peu de particules (moins de 10), mais elle devient moins efficace avec un grand nombre de particules à cause des répulsions mutuelles qui dispersent les particules malgré la convergence initiale.

En résumé, notre méthode \textbf{fonctionne dans certains cas} (46\% des simulations) mais n'est pas fiable de manière générale. Pour de meilleurs résultats, il faudrait une méthode qui tient compte des répulsions pour choisir le placement initial des particules.

\end{document}
